% Appendix A: the formal names (loop0009 item 01). Every Lean-name footnote of
% the draft moved here; the text cites an entry as (Appendix A, A.n) with
% the leanref macro. check_lean.py and check_axioms.py read the leannames
% cells. Axiom line printed by paper1/axiom_check.lean (`make axioms`).
\section{Formal names}
\label{app:lean}

Every theorem, lemma and counterexample of this paper is checked in the Lean~4
proof assistant with the Mathlib library, unless the text says otherwise. The
reader need not read Lean to follow the paper. This appendix is for a reader
who wants to find the formal statement of a result: Table~\ref{tab:lean} maps
each result, in the order in which it appears, to the names of the Lean
declarations that state it. Where a row cites a definition rather than a
theorem, it is the definition the text uses. Lean names use \texttt{repaired}
for what the text calls \emph{corrected}.

All the theorems named depend only on the three standard axioms of Lean's
classical logic, \texttt{propext}, \texttt{Classical.choice} and
\texttt{Quot.sound}; an axiom check in the repository prints this for every
theorem named in the paper. The development has no \texttt{sorry}, Lean's marker for a statement left
unproved.

Two Lean proofs take a route that the text does not describe. The proof of
Theorem~\ref{thm:decides} goes through narrowness \citep{LY2002ref11}, which
counts the closing order forwards rather than reversing it. The proof of
Theorem~\ref{thm:matching} reduces the deficiency form of Hall's theorem to
Mathlib's Hall theorem by adding $\delta$ dummy stacks. One Lean proof takes a
different route from the published one: the proof
that node search is monotone (Theorem~\ref{thm:core}(5)) follows the crusade
argument of \citet{BienstockSeymour1991}, not the reduction to edge search of
\citet{LaPaugh1993} on which \citet{LY2002ref10} rely. Three kinds of fact
are computed rather than proved. The values of $K_2$, $K_{1,3}$ and $K_{3,3}$
after Theorem~\ref{thm:bands}, which show that every end of the two bands
occurs, are computed by brute force from each quantity's own definition. The
search for smaller counterexamples (Appendix~\ref{app:minimality}) and the
results of Section~\ref{sec:results} are computations.

{\small
\begin{longtable}{@{}>{\raggedright}p{1.0cm}>{\raggedright}p{6.6cm}>{\raggedright\arraybackslash}p{8.0cm}@{}}
\caption{Every result of the paper has a formal statement in Lean, listed
here in the order in which the results appear.}\label{tab:lean}\\
\toprule
Entry & Result in the paper & Lean names \\
\midrule
\endfirsthead
\toprule
Entry & Result in the paper & Lean names \\
\midrule
\endhead
\bottomrule
\endfoot
\leangroup{Section~\ref{sec:equiv}, the equivalences}
\leanentry{clique}{Lemma~\ref{lem:clique}, cliques sit in a bag} \leannames{PathDecomposition.exists_bag_of_isClique} \\
\leanentry{vs}{Theorem~\ref{thm:vs}, vertex separation is pathwidth} \leannames{vertexSeparation_eq_pathwidth} \\
\leanentry{mosp}{Theorem~\ref{thm:mosp}, MOSP is pathwidth plus one} \leannames{MOSPInstance.mospValue_eq_pathwidth_add_one} \\
\leanentry{star}{The pattern graph gives no bound: the star $K_{1,4}$ (after Theorem~\ref{thm:mosp})} \leannames{star_mospValue, star_pathwidth_agreementGraph, star_refutes_pattern_graph_bound} \\
\leanentry{core-gml}{Theorem~\ref{thm:core}(1), gate matrix layout} \leannames{NetGateMatrix.tracks_eq_mospValue, NetGateMatrix.tracks_eq_pathwidth_add_one} \\
\leanentry{core-it}{Theorem~\ref{thm:core}(2), interval thickness} \leannames{intervalThickness_eq_pathwidth_add_one} \\
\leanentry{core-logic}{Theorem~\ref{thm:core}(3), one-dimensional logic} \leannames{LogicArray.tracks_eq_intervalThickness, LogicArray.tracks_eq_pathwidth_add_one} \\
\leanentry{core-narrow}{Theorem~\ref{thm:core}(4), narrowness} \leannames{narrowness_eq_pathwidth_add_one} \\
\leanentry{core-ns}{Theorem~\ref{thm:core}(5), node search} \leannames{nodeSearchMonotonicity, nodeSearch_eq_vertexSeparation_add_one, nodeSearch_chain} \\
\leanentry{core-fold}{Theorem~\ref{thm:core}(6), multiple folding} \leannames{NetGateMatrix.foldTracks_eq_pathwidth_add_one} \\
\leanentry{core-vsg}{Theorem~\ref{thm:core}(7), Lengauer's vertex separator game} \leannames{vsg_eq_max} \\
\leanentry{bands}{Theorem~\ref{thm:bands}, the bands} \leannames{pathwidth_le_splitBandwidth_le_pathwidth_add_one, vertexSeparation_le_edgeSearch_le_add_two, vertexSeparation_le_progressiveEdgeSearch_le_add_two} \\
\leanentry{false-pla}{Theorem~\ref{thm:false}(1), PLA folding} \leannames{plaTracks_idMatrix_five, plaTracks_idMatrix_unbounded, plaTracks_pathMatrix_unbounded} \\
\leanentry{false-cw}{Theorem~\ref{thm:false}(2), edge separation read as cutwidth} \leannames{cutwidth_unbounded, modCutwidth_unbounded} \\
\leangroup{Section~\ref{sec:search}, the search}
\leanentry{defs}{Step cost, $\Ext_k$ and $\Sol_k$ (Section~\ref{sec:search-def})} \leannames{stepCost, Solvable, SearchSol} \\
\leanentry{F}{Lemma~\ref{lem:F}, free moves never hurt} \leannames{solvable_iff_searchSol_cl} \\
\leanentry{decides}{Theorem~\ref{thm:decides}, the search decides pathwidth} \leannames{searchSol_empty_iff_pathwidth_add_one_le, searchSol_mospGraph_iff_mospValue_le} \\
\leanentry{dict1}{Table~\ref{tab:dictionary}: closing order from $\emptyset$} \leannames{orderOfLayout} \\
\leanentry{dict2}{Table~\ref{tab:dictionary}: open stacks $b(T)$} \leannames{openStacks_eq_card_boundary} \\
\leanentry{dict3}{Table~\ref{tab:dictionary}: $\cost(T,c)$} \leannames{stepCost_eq_card_boundary_add_one} \\
\leanentry{dict4}{Table~\ref{tab:dictionary}: $\Sol_k(\emptyset)$} \leannames{searchSol_empty_iff_pathwidth_add_one_le} \\
\leanentry{dict5}{Table~\ref{tab:dictionary}: free closure $\cl(T)$} \leannames{mem_cl_iff} \\
\leanentry{dict6}{Table~\ref{tab:dictionary}: definite move, corrected} \leannames{isHereditarilyDefinite_iff_isCommittable} \\
\leanentry{premise}{Lemma~\ref{lem:premise}, what the premise says} \leannames{closeCount_eq, isDefinite_iff} \\
\leanentry{cex}{Counterexample~\ref{cex:definite}} \leannames{definiteMove_counterexample, cexFamily_closed} \\
\leanentry{thm1false}{Chu and Stuckey's Theorem~1 (the definite move), under both readings, and Chu's Theorem~6.3.6 are false} \leannames{chuStuckey_theorem1_false, chuStuckey_theorem1_false_literal, chuThesis_theorem636_false} \\
\leanentry{fink}{Fink's restatement of the definite move is false} \leannames{fink_theorem1_false} \\
\leanentry{hered}{The hereditarily definite premise, and its failure in Counterexample~\ref{cex:definite}} \leannames{IsHereditarilyDefinite, not_isHereditarilyDefinite_cex} \\
\leanentry{submod}{Lemma~\ref{lem:submod}, open stacks are submodular} \leannames{openStacks_submodular} \\
\leanentry{repaired}{Theorem~\ref{thm:repaired}, the corrected definite move is sound} \leannames{solvable_cl_insert_of_hereditarilyDefinite, searchSol_cl_insert_of_hereditarilyDefinite} \\
\leanentry{matching}{Theorem~\ref{thm:matching}, the corrected premise is a matching condition} \leannames{isHereditarilyDefinite_iff_hasDefiniteMatching} \\
\leanentry{tamaki}{Lemma~1 of \citet{Kitsunai2016}, transcribed for digraphs in their terms and proved by their argument; on the symmetric digraph with the extension fixed to the child, its commitment condition is the hereditary premise, our theorem follows from it under its two feasibility hypotheses, it follows from ours, and a five-vertex graph shows those hypotheses are not redundant} \leannames{Tamaki.lemma1, Tamaki.inDeg_submodular, Tamaki.isCommittableExt_symm_iff_hereditarilyDefinite, Tamaki.solvable_cl_insert_of_lemma1, Tamaki.lemma1_symm_of_isCommittable, Tamaki.premise_gap} \\
\leanentry{open1}{The published definite move is right when $\open(q,S)\le1$} \leannames{isHereditarilyDefinite_of_openCount_le_one} \\
\leanentry{subset}{The subset rule is sound} \leannames{searchSol_cl_insert_of_newlyOpened_subset, subsetFilter_sound} \\
\leanentry{tiebreak}{The subset rule needs its tie-break (7-customer graph)} \leannames{noTieBreak_counterexample} \\
\leanentry{isbetter}{The better move's premise, restated} \leannames{isBetter_iff} \\
\leanentry{bettercex}{Chu and Stuckey's Theorem~2 (the better move) is false} \leannames{betterMove_counterexample, chuStuckey_theorem2_false} \\
\leanentry{thm638}{Chu's Theorem~6.3.8 is false} \leannames{chuThesis_theorem638_false, chuThesis_theorem638_false_literal} \\
\leanentry{repbetter}{The corrected better move is sound} \leannames{IsRepairedBetter, searchSol_cl_insert_of_repairedBetter} \\
\leanentry{bugs}{The better move's two requirements: the close count and the rule order} \leannames{bugA_counterexample, bugB_counterexample} \\
\leanentry{oldmove}{The old move is sound} \leannames{searchSol_reinsert, not_searchSol_of_oldMove} \\
\leanentry{reinsert}{The old move needs its reinsertion test (path on five vertices)} \leannames{reinsert_needs_test} \\
\leanentry{memo}{Solvability of a state does not depend on the path to it (the memo)} \leannames{solvable_iff_of_cl_eq} \\
\leanentry{fakeold}{A node with an unrefuted old move can record a solvable state} \leannames{exec_fake_oldMove} \\
\leanentry{nodesound}{Node-soundness of a filter (Section~\ref{sec:theorem})} \leannames{NodeSoundAt} \\
\leanentry{main}{Theorem~\ref{thm:main}, our corrected search is sound} \leannames{exec_repairedFullFilter_sound, exec_repairedFullFilter_pathwidth, exec_repairedFullFilter_mospValue} \\
\leanentry{notsound}{Under the previous algorithm the filter is not node-sound} \leannames{not_codeFilterSound_cexGraph} \\
\leanentry{wrong}{Figure~\ref{fig:wrong}: the previous algorithm refutes 6 and finds 7 on the 34-customer instance, whose optimum is 6 and pathwidth 5; its run is a run of the search with the published filter, which is therefore not sound there} \leannames{previousAlgorithm_wrong_optimum, not_filterSound_codeFilter} \\
\leanentry{minor}{Theorem~\ref{thm:minor}: contracting an edge (merging $v$ into $u$, as the solvers do) or deleting a vertex or an edge never increases pathwidth; along a legal contraction sequence, a minor of pathwidth above $w$ forces pathwidth above $w$, and on the MOSP graph more than $w+1$ open stacks} \leannames{pathwidth_contractEdge_le, pathwidth_induce_le, pathwidth_mono, pathwidth_minorOf_le, pathwidth_gt_of_minorOf, MOSPInstance.lt_mospValue_of_minorOf} \\
\leanentry{grid}{The $(h,w)$-grid has pathwidth $\min(h,w)$, except the $1\times1$ grid (pathwidth $0$); \citet[Theorem~4.1]{EllisWarren2008}. Upper bound by the row-major layout's bandwidth; lower bound directly on vertex separation, at the first prefix that completes a row or a column} \leannames{pathwidth_grid, pathwidth_grid_le, min_le_pathwidth_grid, pathwidth_grid_one_one, pathwidth_grid_square} \\
\leanentry{hypercube}{The $d$-dimensional hypercube has pathwidth $\sum_{i<d}\binom{i}{\lfloor i/2\rfloor}$ \citep[Theorem~2]{ChandranKavitha2006}, with Harper's vertex-isoperimetric theorem \citep{Harper1966} proved by compression. Upper bound by the layout along Harper's simplicial order, lower bound at a prefix of every size, both by induction on $d$ and on vertex separation} \leannames{pathwidth_hypercube, pathwidth_cubeGraph, Hypercube.harper, Hypercube.card_nbhd_init_le, Hypercube.exists_init_profile} \\
\leanentry{mycielski}{For every graph $G$, its Mycielskian $\mu(G)$ (originals, a shadow of each vertex adjacent to its neighbours, and an apex adjacent to the shadows) has $\mathrm{pw}(G)+1\le\mathrm{pw}(\mu(G))\le 2\,\mathrm{pw}(G)+1$. The upper bound interleaves each vertex with its shadow and places the apex last; we found no published statement of it. The lower bound keeps the bags meeting the apex and the shadows} \leannames{pathwidth_mycielskian_le, pathwidth_add_one_le_mycielskian} \\
\leanentry{emptypair}{Step~1 of Appendix~\ref{app:random}: at every position of every layout the prefix and the inactive suffix are an empty pair, so for every $j\le n$ some $j$-set has closed neighbourhood of size at most $j+\pw$ (Harper's bound in layout form), and $n-2s^{*}-2\le\pw$ for the empty-pair index $s^{*}$} \leannames{exists_empty_pair_of_layout, exists_closedNbhd_card_le, exists_empty_pair, card_sub_two_mul_emptyPairIndex_sub_two_le_pathwidth} \\
\leanentry{firstmoment}{Step~2 of Appendix~\ref{app:random}: the number of graphs on $n$ labelled vertices of pathwidth at most $k$ is at most $\binom nt\binom{n-t}{t-1}2^{\binom n2-t(t-1)}$, $t=\lceil(n-k)/2\rceil$, and the $G(n,p)$-weighted sum of those graphs is at most $\binom nt\binom{n-t}{t-1}(1-p)^{t(t-1)}$; there are $2^{\binom n2}$ graphs on $n$ labelled vertices} \leannames{card_pathwidth_le_le, weightedSum_graphsOfPathwidthLe_le, card_emptyPairGraphs, card_simpleGraph} \\
\leanentry{firstmomentA}{Theorem~\ref{thm:firstmoment} and its consequence: the weighted sum is at most $\exp(-\delta n+3\ln n+ca)$ for $n\ge\max(3,c)$, at most $e^{c/2}n^{3}e^{-\delta n}$, and tends to $0$ when $\delta>0$; the trinomial entropy bound from the binomial theorem} \leannames{firstMoment_pathwidth_le_fin, firstMoment_pathwidth_le_fin', tendsto_firstMoment_pathwidth, choose_mul_choose_le_exp_entropy} \\
\leanentry{clique-cover}{Appendix~\ref{app:random-cs}: for independent uniform patterns of given sizes, the probability that the instance's MOSP value is at most $k+1$ (its pathwidth at most $k$) is at most $\binom nt\binom{n-t}{t-1}\prod_j N(s_j)/\binom n{s_j}$, $N(s)$ the number of $s$-sets missing $A$ or missing $R$} \leannames{prob_coversOfMospValueLe_le, prob_coversOfPathwidthLe_le, card_coversOfPathwidthLe_le, card_filter_avoid} \\
\leanentry{coupling}{Appendix~\ref{app:random-cs}, the coupling: if every pattern of one cover is contained in the corresponding pattern of another, every empty pair of the larger cover's graph is an empty pair of the smaller's, and the pathwidth is monotone} \leannames{isEmptyPair_coverGraph_of_subset, pathwidth_coverGraph_mono} \\
\leanentry{blowup}{Theorem~\ref{thm:blowup}, the blow-up bound: for every instance, $\mathrm{MOSP}\le\omega\,(\pw(B)+1)$ with $\omega$ the largest pattern and $B$ the customer--pattern incidence graph (as a \texttt{SimpleGraph} on the sum type); a path decomposition of $B$ blown up to one of the MOSP graph} \leannames{MOSPInstance.mospValue_le_maxPatternSize_mul, MOSPInstance.pathwidth_mospGraph_add_one_le, MOSPInstance.blowUp_width_add_one_le} \\
\leanentry{forest}{Theorem~\ref{thm:blowup}, the forest form: if $B$ is acyclic, $\mathrm{MOSP}\le\omega(\lfloor\log_3((6(n+m)+3)/5)\rfloor+1)$; a forest on $N$ vertices has pathwidth at most $\lfloor\log_3((6N+3)/5)\rfloor$, the least $k$ with $N\le(5\cdot3^k-3)/2$; a tree on at most $(5\cdot3^k-3)/2$ vertices has pathwidth at most $k$} \leannames{MOSPInstance.mospValue_le_of_incidence_isAcyclic, pathwidth_le_treeBoundInv_of_isAcyclic, le_treeBound_treeBoundInv, treeBound_lt_of_lt_treeBoundInv, tree_pathwidth_le_treeBound} \\
\leanentry{mono}{Appendix~\ref{app:excess}: the MOSP value is monotone under adding requirements (Yanasse's subinstance monotonicity)} \leannames{MOSPInstance.mospValue_mono} \\
\leanentry{layer}{Lemma~\ref{lem:layer}, the layer lemma: for any level function with adjacent vertices at most one level apart, $\pw\le\max_i(|L_i|+|L_{i+1}|)-1$, in pathwidth and vertex-separation form, and for the distance layers from any vertex of a connected graph} \leannames{pathwidth_le_layerWidth, vertexSeparation_le_layerWidth, pathwidth_le_layerWidth_dist} \\
\leanentry{indep}{Lemma~\ref{lem:layer}, the independent-set bound $\pw\le n-\alpha$, by the bags $(V\setminus I)\cup\{v\}$} \leannames{pathwidth_le_card_sub_indepNum} \\
\leanentry{pendant}{Lemma~\ref{lem:pendant}, the two-sided pendant bound: a clique $R$ with $q\ge1$ cliques of size at least $t$ meeting $R$ in one vertex each, at distinct vertices, forces $\pw\ge t+\lceil q/2\rceil-2$; for any graph, for one decomposition, in the depth-one case $t=q=|R|$, and in MOSP terms (a pattern with $q$ patterns of size at least $t$ hanging off distinct customers forces $\mathrm{MOSP}\ge t+\lceil q/2\rceil-1$)} \leannames{pathwidth_ge_of_pendant_cliques, PathDecomposition.pendant_bound, pathwidth_ge_of_clique_pendant_cliques, MOSPInstance.mospValue_ge_of_pendant_patterns} \\
\leanentry{separator}{Theorem~\ref{thm:separator}, the deterministic step: a layout of width $k<n$ forces pairwise disjoint $A,R,S$ covering $V$ with $|A|=\lceil(n-k)/2\rceil$, $|R|=\lfloor(n-k)/2\rfloor$, $|S|=k$ and no $A$--$R$ edge (on the MOSP graph: no pattern meets both $A$ and $R$); a walk avoiding $S$ that starts in $A$ stays in $A$, so every component of $G-S$ has at most $\lceil(n-k)/2\rceil$ vertices; $3K_2$ at $k=0$ satisfies the component condition and not the bipartition one} \leannames{exists_balanced_separator, exists_balanced_separator_cover, exists_isBalancedBip, componentOff_subset_side, card_componentOff_le_of_isBalancedBip, not_isBalancedBip_threeK2, card_componentOff_threeK2_le} \\
\leanentry{closedpair}{The closed pair (Appendix~\ref{app:excess}): $R\subseteq V\setminus N[A]$ iff $(A,R)$ is disjoint and empty; one round of the alternating closure $R^{*}=V\setminus N[A]$, $A^{*}=V\setminus N[R^{*}]$ is monotone on empty pairs, lands on a closed pair from every $A$, is idempotent, has separator $N(A^{*})\cap N(R^{*})$ and keeps the sizes; a pair is closed iff it is empty and every outside vertex sees both sides; a layout of width $k<n$ forces a closed empty pair with sides of at least the forced sizes} \leannames{subset_compl_closedNbhd_iff, subset_closeStep, closeStep_isClosedPair, closeStep_closeStep, sdiff_eq_inter_boundary_of_isClosedPair, card_closeStep, isClosedPair_iff_sees_both, exists_closedPair_of_pathwidth_le} \\
\leanentry{weightedcount}{The weighted count (Appendix~\ref{app:excess}): for any weight $1\le Z(P)\le\#\{\text{empty pairs}\}$ on the empty pairs, $\sum_P \mathbf 1[P\ \text{empty}]/Z(P)\ge\mathbf 1[\text{some pair is empty}]$, so the $G(n,p)$-weighted sum of the graphs of pathwidth at most $k$ is at most $E[\sum_P \mathbf 1[P\ \text{empty}]/Z(G,P)]$ over the exact-size pairs; with $Z$ the number of flips of $P$ (exchanging equally many vertices isolated in $G[A]$ and $G[R]$ gives distinct empty pairs of the same sizes), the flip-weighted count} \leannames{weightedCount_ge_indicator, weightedSum_graphsOfPathwidthLe_le_weighted, flipPair_isEmptyPair, flipPair_injOn, card_flipPair, card_flips_le, weightedSum_graphsOfPathwidthLe_le_flipWeighted} \\
\leanentry{theoremS}{Theorem~\ref{thm:separator}: for $0\le p\le1$ and $k<n$, the $G(n,p)$-weighted sum of the graphs of pathwidth at most $k$ is at most $\binom nk$ times the $G(N,p)$-weighted count, $N=n-k$, of the graphs on $V\setminus S$ whose vertices split into an empty pair of sizes $\lceil N/2\rceil$, $\lfloor N/2\rfloor$; the union bound over the $\binom nk$ separators is the only inequality: the graphs on $V$ with a given restriction to $V\setminus S$ have total weight equal to the restriction's weight (the pairs meeting $S$ carry weight $(p+(1-p))^{|\text{pairs meeting }S|}=1$), and the inner count is the same for every $S$ by a relabelling of $V$} \leannames{weightedSum_graphsOfPathwidthLe_le_choose_mul_innerBip, weightedSum_graphsOfPathwidthLe_le_separator, weightedSum_graphsOfPathwidthLe_le_sum_separators, sum_gnpWeight_fiber, sum_gnpWeight_bipGraphs, card_pairsOff, innerBip_eq_of_card_eq} \\
\leanentry{theoremCl}{The closed-pair count (Appendix~\ref{app:excess}), the union bound: the weighted sum of the graphs of pathwidth at most $k$ is at most $\sum_{(A,R)}\Pr[(A,R)\ \text{is closed}]$ over the ordered disjoint pairs with $|A|\ge\lceil(n-k)/2\rceil$, $|R|\ge\lfloor(n-k)/2\rfloor$; each term evaluated in the next row} \leannames{weightedSum_graphsOfPathwidthLe_le_sum_closedPairs} \\
\leanentry{theoremClproduct}{The closed-pair count, the independence step: the $p$-weight of the subsets of a set of pairs that satisfy one predicate on each of pairwise disjoint blocks is the product of the per-block weights (the pairs in no block are free), so the $G(n,p)$ probability that $(A,R)$ is closed is exactly $(1-p)^{|A||R|}\bigl((1-(1-p)^{|A|})(1-(1-p)^{|R|})\bigr)^{|S|}$ (one block of the forbidden $A$--$R$ pairs, two star blocks per $v\in S$); hence $\Pr[\pw\le k]\le\sum_{a\ge t,\,r'\ge r}\binom na\binom{n-a}{r'}(1-p)^{ar'}\bigl((1-(1-p)^{a})(1-(1-p)^{r'})\bigr)^{n-a-r'}$, the expected number of closed pairs of at least the forced sizes} \leannames{sum_subsetWeight_split, sum_subsetWeight_blocks, sum_gnpWeight_closedPairGraphs, weightedSum_graphsOfPathwidthLe_le_closedPairCount, weightedSum_graphsOfPathwidthLe_le_closedPairFormula} \\
\leanentry{reduction-theorem}{Theorem~3.2 of \citet{EdwardsMcDermid2015}, the reduction theorem
  (Appendix~\ref{app:random}): every graph with every degree $0$ or at least $3$ has a vertex set $X$
  such that $G-X$ reduces to a non-empty basic graph $B$ with $|X|+f^-(B)\le f^-(G)$; unconditional,
  with the step inequality in all four cases and a local commutation lemma in place of the appeal to
  the uniqueness of the reduced graph} \leannames{exists_reduction_set, one_add_potentialMinus_le, SPReduces.isolateSet_comm, potential_isolate} \\
\leanentry{degree-bound}{Lemma~5.1 and Corollary~5.2 of \citet{EdwardsMcDermid2015}
  (Appendix~\ref{app:random}): $\pw(G)\le\sum_v g(d_G(v))+\tau(n)+3\rho(n)+1$ for every $G$,
  conditional on two cited results stated as named propositions and never asserted --- the cubic
  base case \path{CubicPathwidthBound} with budget $\tau$ \citep[Theorem~5]{FominHoie2006} and the
  reduction lemma \path{ReductionPathwidthBound} with budget $\rho$
  \citep{KneisMolleRichterRossmanith2009} --- and on $\tau$ monotone and non-negative; the paper
  gives no constants} \leannames{pathwidth_le_histogram_of_hypotheses, pathwidth_le_degreeSum_of_hypotheses, pathwidth_le_potentialMinus_of_hypotheses, pathwidth_basic_le, CubicPathwidthBound, ReductionPathwidthBound} \\
\leanentry{endpoint}{Lemma~\ref{lem:endpoint}, the premises as commitment conditions} \leannames{isDefinite_iff_endpoint, isHereditarilyDefinite_iff_isCommittable} \\
\leanentry{cexnotcommit}{In Counterexample~\ref{cex:definite} the endpoint test passes and the commitment fails} \leannames{cex_isDefinite_not_isCommittable} \\
\leanentry{existscommit}{The published premise gives a commitment to the least-border set} \leannames{exists_isCommittable_of_isDefinite} \\
\leanentry{commit234}{In Counterexample~\ref{cex:definite} that set is $\{2,3,4\}$} \leannames{cex_isCommittable_234} \\
\leanentry{depth1}{A commitment of depth one holds iff $\open(q,S)\le1$} \leannames{isCommittable_insert_iff} \\
\leanentry{greedy}{The greedy step of Coudert, Mazauric and Nisse is that condition} \leannames{isGreedyStep_iff} \\
\leangroup{Section~\ref{sec:results}, computational results}
\leanentry{split}{The split of the root into independent tasks is sound} \leannames{execSplit_repairedFullFilter_sound, execSplit_repairedFullFilter_mospValue} \\
\leangroup{Appendix~\ref{app:machinery-all}, supporting material}
\leanentry{edgeless-ns}{Node search equalities need an edge (Appendix~\ref{app:edge})} \leannames{nodeSearch_ne_intervalThickness_of_edgeless, monotoneNodeSearch_ne_vertexSeparation_add_one_of_edgeless} \\
\leanentry{vsg0}{Lengauer's Theorem~4 fails at $k=0$ (Appendix~\ref{app:edge})} \leannames{isPositiveVSG_triangleGraph_counterexample} \\
\leanentry{tns}{$t=\ns$ fails on the matrix $[1]$ (Appendix~\ref{app:edge})} \leannames{tracks_ne_monotoneNodeSearch_one} \\
\leanentry{kp}{The Kirousis--Papadimitriou claim fails on $K_{1,3}$ (Appendix~\ref{app:edge})} \leannames{kirousisPapadimitriou_claim2_false} \\
\leanentry{pebbling}{Progressive pebbling is in the exact core (Appendix~\ref{app:thirteenth})} \leannames{pbw_lengauerD_eq_pathwidth, pbw_eq_mospValue_pebbleMatrix, mpb_eq_pathwidth_add_one} \\
\leanentry{vsearch}{The vertex search is in the exact core (Appendix~\ref{app:thirteenth})} \leannames{vertexSearchNumber_eq_pathwidth_add_one, monotoneVertexSearchNumber_eq_pathwidth_add_one, isNodeSearch_of_isSuccessfulVertexSearch, vertexSearchNumber_eq_nodeSearch, vertexSearchNumber_of_isEmpty} \\
\leanentry{wavefront}{The maximum wavefront is in the exact core (Appendix~\ref{app:thirteenth})} \leannames{card_wavefront, maxWavefront_eq_vertexSepOfLayout_add_one, minMaxWavefront_eq_pathwidth_add_one, minMaxWavefront_eq_narrowness, minMaxWavefront_of_isEmpty} \\
\leanentry{filter}{Proposition~\ref{prop:filter}, the corrected filter is node-sound} \leannames{repairedFullFilter_sound} \\
\leanentry{runs}{Theorem~\ref{thm:runs}, runs are sound} \leannames{Exec.sound} \\
\leanentry{encoding}{Appendix~\ref{app:sat}, the SAT formulation is faithful (clauses 1--4)} \leannames{encodes_iff_mospValue_le} \\
\leanentry{obstruction}{A checked lower-bound certificate implies the bound (Section~\ref{sec:certs}): an UNSAT expansion level $(t,b)$ with $t\le n$ gives $b+1\le t+\pw$ (Harper's bound, contrapositive; the hypothesis $t\le n$ is necessary, one-vertex counterexample); a legal contraction sequence whose minor has minimum degree $d$ gives $d\le\pw$, in the \texttt{minDegreeIn} form and in the degree form on a nonempty minor (necessary, likewise); the \texttt{Bool} replay \path{checkSeq} is sound and complete for legal sequences; both in MOSP units} \leannames{pathwidth_ge_of_forall_closedNbhd, pathwidth_ge_of_unsat, pathwidth_ge_of_forall_closedNbhd_needs_le, mospValue_ge_of_unsat, checkSeq_sound, checkSeq_complete, checkSeq_iff, pathwidth_ge_of_minor_minDegree, pathwidth_ge_of_checkSeq, degree_minorOf_eq, pathwidth_ge_of_minor_forall_degree, pathwidth_ge_of_minor_forall_degree_needs_nonempty, mospValue_ge_of_minor_minDegree} \\
\leanentry{obstruction-decided}{Three certificates closed by \texttt{decide}: the 8-cycle contracted to a triangle ($2\le\pw$), $K_{3,4}$ at level one over all 128 subsets ($3\le\pw$), and the record instance \emph{Warwick 1729} (30 customers, 90 requirements), whose 22-step contraction sequence to $K_8$ is replayed on the MOSP graph Lean builds from the matrix: every pattern order opens at least 8 stacks (optimum 9)} \leannames{two_le_pathwidth_octGraph, three_le_pathwidth_k34, seven_le_pathwidth_warwick, eight_le_mospValue_warwick} \\
\end{longtable}
}
